\documentclass[12pt]{article}
\usepackage[utf8]{inputenc}
\usepackage[T1]{fontenc}
\usepackage{hyperref}
\title{Settla Language Resources for K-12 Curriculum Development}
\author{Compiled by Claude Sonnet 5.5 (claude-sonnet-5-5) for Andr\'as Kornai}
\date{}
\begin{document}
\maketitle

\section{Introduction}
Settla (also called Kisetla or Settler Swahili; ISO 639-3: sta) is a Swahili-based pidgin with no native speakers and no ethnic community. Sources describe it as used by European settlers to talk with Swahili speakers. The Endangered Languages Project places it in Zambia's Central province, north of Lusaka. Verbix classifies it as dormant, and Wikipedia says a thorough study of its morphology has not been done. The search found no educational material for Settla, and none is claimed here.

\section{Primary Educational Resources}
No primers, textbooks or school materials were found. A pidgin with no first-language speakers and no community is not a realistic base for a K-12 curriculum in the usual sense.

\section{Linguistic and Reference Materials}
\begin{itemize}
\item Wikipedia article on Settler Swahili: \url{https://en.wikipedia.org/wiki/Settler_Swahili}
\item Endangered Languages Project page for Settla (a bibliography page also appeared in the results, which I did not open): \url{https://endangeredlanguages.com/lang/5419}
\item Verbix page on Settla, with limited vocabulary and grammar notes: \url{https://docs.verbix.com/EndangeredLanguages/Settla}
\item Wikipedia article on the languages of Zambia, which lists Settla with Fanagalo among Zambia's pidgins: \url{https://www.wikipedia.com/wiki/Zambian_language}
\end{itemize}

\section{Religious and Community Resources}
None found.

\section{Digital and Audio Resources}
None found.

\section{Curriculum Development Context}
Settla is used only as a second language and has no ethnic community. Little is documented, and the sources disagree on whether it was mainly a Zambian or a Kenyan phenomenon. Realistic options are historical or sociolinguistic documentation, or teaching about it inside Swahili or English language courses. A curriculum that uses Settla as the medium of instruction does not look feasible from what was found.

\end{document}